\begin{proof}[Proof of \cref{obj:thm:rare-label-floor}]
\begin{enumerate}
\item \textit{Choice of the two laws.}
Fix \(n,m,d,\epsilon\) satisfying the hypotheses of the risk bounds, and set
\[
c:=\frac1{1024},
\qquad
S:=n\epsilon>0,
\qquad
\delta_n:=\frac1{16\sqrt{\max\{1,S\}}}.
\]
Since \(\max\{1,S\}\ge1\), we have
\[
0<\delta_n\le\frac1{16}\le\frac18.
\]
Write
\[
\mathsf Z_d:=\{1,\ldots,d\}\times\{0,1\}^2.
\]
For \(\xi\in\{-1,1\}\), define the law \(P_{\xi,n}\) on \(\mathsf Z_d\) by
\[
\begin{aligned}
P_{\xi,n}(j,1,1)&=\frac{\epsilon}{2}\left(\frac12+\xi\delta_n\right),&
P_{\xi,n}(j,1,0)&=\frac{\epsilon}{2}\left(\frac12-\xi\delta_n\right),\\
P_{\xi,n}(j,0,1)&=\frac{1-\epsilon}{4},&
P_{\xi,n}(j,0,0)&=\frac{1-\epsilon}{4},
&&j=1,2,
\end{aligned}
\]
and
\[
P_{\xi,n}(j,a,y)=0,
\qquad j>2,\quad a,y\in\{0,1\}.
\]
Also set
\[
P_{+,n}:=P_{1,n},
\qquad
P_{-,n}:=P_{-1,n}.
\]
Every atom in the first two covariate cells is positive. Summing within each such cell gives
\[
\sum_{a,y\in\{0,1\}}P_{\xi,n}(j,a,y)
=
\frac{\epsilon}{2}+\frac{1-\epsilon}{2}
=\frac12,
\qquad j=1,2.
\]
Thus the total mass is one. Summing over outcomes and taking the conditional ratios gives
\[
\begin{aligned}
p_j(P_{\xi,n})&=\frac12,
&
e_j(P_{\xi,n})&=\frac{\epsilon/2}{1/2}=\epsilon,\\
\mu_{0j}(P_{\xi,n})&=
\frac{(1-\epsilon)/4}{(1-\epsilon)/2}=\frac12,
&
\mu_{1j}(P_{\xi,n})&=
\frac{(\epsilon/2)(1/2+\xi\delta_n)}{\epsilon/2}
=\frac12+\xi\delta_n,
&&j=1,2.
\end{aligned}
\]
The remaining covariate masses are zero. Because
\[
\epsilon\le1-\epsilon,
\]
the occupied cells satisfy \cref{obj:ass:overlap}, so \cref{obj:def:model} gives
\[
P_{+,n},P_{-,n}\in\mathcal M_{d,\epsilon}.
\]
By \cref{obj:def:target},
\[
\tau(P_{\xi,n})
=
\sum_{j=1}^{2}\frac12\,\xi\delta_n
=\xi\delta_n.
\]
Define their common treatment--covariate marginal by
\[
\pi_{XA}(j,a):=
\begin{cases}
\epsilon/2,&j\in\{1,2\},\ a=1,\\
(1-\epsilon)/2,&j\in\{1,2\},\ a=0,\\
0,&j>2.
\end{cases}
\]
The outcome sums above show that
\[
(P_{+,n})_{XA}=(P_{-,n})_{XA}=\pi_{XA}.
\]
When \(d=2\), all statements concerning \(j>2\) have an empty index set.
% lean: CausalSmith.Stat.AnnotationRarearmFrontier.labelFloorAmplitude_scaled
% lean: CausalSmith.Stat.AnnotationRarearmFrontier.labelFloor_jointMass
% lean: CausalSmith.Stat.AnnotationRarearmFrontier.labelFloor_model
% lean: CausalSmith.Stat.AnnotationRarearmFrontier.labelFloor_target
% lean: CausalSmith.Stat.AnnotationRarearmFrontier.labelFloor_auxMarginal

\item \textit{Closeness of the complete-record experiments.}
Use the natural-logarithm divergence and finite total variation distance of
\cref{obj:lem:pinsker-finite}. The two laws have common positive support,
namely \(\{1,2\}\times\{0,1\}^2\), so their divergence is finite.
For either occupied cell, cancellation of the common control terms gives
\[
\begin{aligned}
&\sum_{a,y\in\{0,1\}}
P_{+,n}(j,a,y)
\log\frac{P_{+,n}(j,a,y)}{P_{-,n}(j,a,y)}\\
&\quad=
\frac{\epsilon}{2}\left[
\left(\frac12+\delta_n\right)
\log\frac{1/2+\delta_n}{1/2-\delta_n}
+
\left(\frac12-\delta_n\right)
\log\frac{1/2-\delta_n}{1/2+\delta_n}
\right]\\
&\quad=
\epsilon\delta_n
\log\frac{1/2+\delta_n}{1/2-\delta_n},
\qquad j=1,2,
\end{aligned}
\]
where the last equality uses \(\log(x^{-1})=-\log x\).
The null cells contribute zero. Hence
\[
D(P_{+,n}\Vert P_{-,n})
=
2\epsilon\delta_n
\log\frac{1/2+\delta_n}{1/2-\delta_n}.
\]
The ratio is positive, and \(\log x\le x-1\) yields
\[
\log\frac{1/2+\delta_n}{1/2-\delta_n}
\le
\frac{1/2+\delta_n}{1/2-\delta_n}-1
=
\frac{2\delta_n}{1/2-\delta_n}
\le8\delta_n.
\]
Indeed, \(1/2-\delta_n\ge3/8\ge1/4\). Multiplication by
\(2\epsilon\delta_n>0\) therefore gives
\[
D(P_{+,n}\Vert P_{-,n})
\le16\epsilon\delta_n^2.
\]

Define the complete-record array laws by
\[
\mathsf P_{\mathrm{lab},\pm}:=P_{\pm,n}^{\otimes n}.
\]
They too have common support. For an array
\(\boldsymbol z=(z_1,\ldots,z_n)\) in that support,
\[
\log
\frac{\mathsf P_{\mathrm{lab},+}(\boldsymbol z)}
     {\mathsf P_{\mathrm{lab},-}(\boldsymbol z)}
=
\sum_{i=1}^{n}
\log\frac{P_{+,n}(z_i)}{P_{-,n}(z_i)}.
\]
Taking expectations, each coordinate has law \(P_{+,n}\), and thus
\[
\mathbb E_{\mathsf P_{\mathrm{lab},+}}
\left[
\sum_{i=1}^{n}
\log\frac{P_{+,n}(z_i)}{P_{-,n}(z_i)}
\right]
=
nD(P_{+,n}\Vert P_{-,n}).
\]
Consequently the product divergence is finite and satisfies
\[
D(\mathsf P_{\mathrm{lab},+}\Vert\mathsf P_{\mathrm{lab},-})
\le
nD(P_{+,n}\Vert P_{-,n})
\le16S\delta_n^2.
\]
Squaring the chosen amplitude gives
\[
\delta_n^2
=
\frac1{256\max\{1,S\}},
\qquad
16S\delta_n^2
=
\frac{S}{16\max\{1,S\}}
\le\frac1{16}.
\]
Applying \cref{obj:lem:pinsker-finite} to these finite array laws now gives
\[
\operatorname{TV}
(\mathsf P_{\mathrm{lab},+},\mathsf P_{\mathrm{lab},-})
\le
\sqrt{
\frac{D(\mathsf P_{\mathrm{lab},+}\Vert
        \mathsf P_{\mathrm{lab},-})}{2}}
\le\frac14,
\]
since the squared expression under the final comparison is at most
\(1/32\le1/16\).
% lean: CausalSmith.Stat.AnnotationRarearmFrontier.labelFloor_KL_formula
% lean: CausalSmith.Stat.AnnotationRarearmFrontier.labelFloor_KL_bound
% lean: CausalSmith.Stat.AnnotationRarearmFrontier.labelFloor_absolutelyContinuous
% lean: CausalSmith.Stat.AnnotationRarearmFrontier.labelFloor_labeled_KL_bound
% lean: CausalSmith.Stat.AnnotationRarearmFrontier.labelFloor_scaled_tv

\item \textit{Calibration to the benchmark.}
The square identity gives
\[
\frac{\delta_n^2}{4}
=
\frac1{1024\max\{1,S\}}.
\]
If \(S\le1\), positivity of \(S\) gives \(S^{-1}\ge1\), and therefore
\[
\frac1{\max\{1,S\}}=1=\min\{1,S^{-1}\}.
\]
If \(S>1\), then \(S^{-1}<1\), and
\[
\frac1{\max\{1,S\}}=S^{-1}=\min\{1,S^{-1}\}.
\]
In both cases,
\[
\frac{\delta_n^2}{4}
=
\frac1{1024}\min\{1,(n\epsilon)^{-1}\}
=
c\,b(n,\epsilon).
\]
% lean: CausalSmith.Stat.AnnotationRarearmFrontier.labelFloor_scaled_rate

\item \textit{Adjoining the auxiliary observations and randomization.}
For probability measures on a general measurable space, use the extension
\[
\operatorname{TV}(\nu,\nu')
:=
\sup_{\mathcal E\ \mathrm{measurable}}
|\nu(\mathcal E)-\nu'(\mathcal E)|,
\]
which agrees with the finite formula in \cref{obj:lem:pinsker-finite}.

We first establish the common-factor identity needed here. Let
\(\nu,\nu'\) be probability laws on a finite set
\(\Omega_{\mathrm{fin}}\), and let \(\rho\) be a probability law on a
measurable space \(\Omega_{\mathrm{common}}\). For any measurable
\(\mathcal E\subseteq
\Omega_{\mathrm{fin}}\times\Omega_{\mathrm{common}}\), define its sections by
\[
\mathcal E_z:=\{w\in\Omega_{\mathrm{common}}:(z,w)\in\mathcal E\},
\qquad z\in\Omega_{\mathrm{fin}}.
\]
The signed masses sum to zero, so
\[
\sum_z\max\{\nu(z)-\nu'(z),0\}
=
\sum_z\max\{\nu'(z)-\nu(z),0\}
=
\frac12\sum_z|\nu(z)-\nu'(z)|.
\]
Since \(0\le\rho(\mathcal E_z)\le1\), it follows that
\[
\begin{aligned}
|(\nu\otimes\rho)(\mathcal E)-(\nu'\otimes\rho)(\mathcal E)|
&=
\left|
\sum_z(\nu(z)-\nu'(z))\rho(\mathcal E_z)
\right|\\
&\le\frac12\sum_z|\nu(z)-\nu'(z)|
=\operatorname{TV}(\nu,\nu').
\end{aligned}
\]
Conversely, for every subset
\(\mathcal E_{\mathrm{fin}}\subseteq\Omega_{\mathrm{fin}}\),
\[
(\nu\otimes\rho)
(\mathcal E_{\mathrm{fin}}\times\Omega_{\mathrm{common}})
=
\nu(\mathcal E_{\mathrm{fin}}),
\]
and the same identity holds for \(\nu'\). Taking suprema in the two
directions proves
\[
\operatorname{TV}(\nu\otimes\rho,\nu'\otimes\rho)
=
\operatorname{TV}(\nu,\nu').
\]

By \cref{obj:ass:data-product}, the augmented observation laws for the
original-record experiment are
\[
\mathsf P_{\mathrm{ann},\pm}
:=
\bigl(\mathsf P_{\mathrm{lab},\pm}
      \otimes\pi_{XA}^{\otimes m}\bigr)
\otimes\mathrm{Uniform}[0,1].
\]
First removing the common seed factor and then the common auxiliary
factor gives
\[
\begin{aligned}
\operatorname{TV}
(\mathsf P_{\mathrm{ann},+},\mathsf P_{\mathrm{ann},-})
&=
\operatorname{TV}
\bigl(
\mathsf P_{\mathrm{lab},+}\otimes\pi_{XA}^{\otimes m},
\mathsf P_{\mathrm{lab},-}\otimes\pi_{XA}^{\otimes m}
\bigr)\\
&=
\operatorname{TV}
(\mathsf P_{\mathrm{lab},+},\mathsf P_{\mathrm{lab},-})
\le\frac14.
\end{aligned}
\]
For \(m=0\), the auxiliary array space consists of the empty array and its
law is a point mass, so the same argument applies.
% lean: CausalSmith.Stat.AnnotationRarearmFrontier.labelFloor_tv_common_product

\item \textit{Bounded risks and the testing reduction.}
For any observable law \(P\), the conditional outcome means satisfy
\begingroup\small
\[
\mu_{aj}(P)=
\begin{cases}
\displaystyle
\frac{P(j,a,1)}{P(j,a,0)+P(j,a,1)},
& P(j,a,0)+P(j,a,1)>0,\\[6pt]
0,&P(j,a,0)+P(j,a,1)=0,
\end{cases}
\qquad a\in\{0,1\},\quad j\in\{1,\ldots,d\}.
\]
\endgroup
The numerator lies between zero and the denominator, so
\[
0\le\mu_{aj}(P)\le1,
\qquad
-1\le\mu_{1j}(P)-\mu_{0j}(P)\le1.
\]
Since \(p_j(P)\ge0\) and \(\sum_jp_j(P)=1\),
\cref{obj:def:target} yields
\[
-1
=
\sum_jp_j(P)(-1)
\le\tau(P)
\le
\sum_jp_j(P)
=1.
\]
Consequently, every clipped rule obeys
\[
|T(\mathcal D,U)-\tau(P)|\le2,
\qquad
0\le(T(\mathcal D,U)-\tau(P))^2\le4,
\]
and hence its risk is at most four under every model law. The same
pointwise bound applies to clipped exact-marginal rules. The constant
zero rule belongs to both rule classes.

To establish the testing inequality, let
\[
\nu_\pm\in\mathsf{Prob}(\Omega_{\mathrm{test}}),
\qquad
T_{\mathrm{test}}:\Omega_{\mathrm{test}}\longrightarrow[-1,1]
\quad\text{be measurable},
\qquad
\operatorname{TV}(\nu_+,\nu_-)\le\frac14,
\]
where \(\Omega_{\mathrm{test}}\) is any measurable space. Define
\[
\mathcal E_+
:=
\{\omega:|T_{\mathrm{test}}(\omega)-\delta_n|\ge\delta_n\},
\qquad
\mathcal E_-
:=
\{\omega:|T_{\mathrm{test}}(\omega)+\delta_n|\ge\delta_n\}.
\]
The targets are separated by
\[
|\delta_n-(-\delta_n)|=2\delta_n.
\]
If \(\omega\in\mathcal E_+^c\), the triangle inequality gives
\[
2\delta_n
\le
|T_{\mathrm{test}}(\omega)-\delta_n|
+
|T_{\mathrm{test}}(\omega)+\delta_n|
<
\delta_n+|T_{\mathrm{test}}(\omega)+\delta_n|.
\]
Thus \(\mathcal E_+^c\subseteq\mathcal E_-\), and the defining event bound
for total variation implies
\[
\begin{aligned}
\nu_+(\mathcal E_+)+\nu_-(\mathcal E_-)
&\ge
\nu_+(\mathcal E_+)+\nu_-(\mathcal E_+^c)\\
&=
1+\nu_+(\mathcal E_+)-\nu_-(\mathcal E_+)\\
&\ge1-\operatorname{TV}(\nu_+,\nu_-).
\end{aligned}
\]
Therefore
\[
\frac14
\le
\frac{1-\operatorname{TV}(\nu_+,\nu_-)}2
\le
\max\{\nu_+(\mathcal E_+),\nu_-(\mathcal E_-)\}.
\]
For each sign, the pointwise loss bound on the corresponding error
region is
\[
\delta_n^2\mathbf1_{\mathcal E_\pm}(\omega)
\le
\bigl(T_{\mathrm{test}}(\omega)\mp\delta_n\bigr)^2.
\]
The losses are integrable because the rule is bounded. Integrating and
taking the maximum yields
\[
\begin{aligned}
\frac{\delta_n^2}{4}
&\le
\delta_n^2
\max\{\nu_+(\mathcal E_+),\nu_-(\mathcal E_-)\}\\
&\le
\max\left\{
\int (T_{\mathrm{test}}-\delta_n)^2\,d\nu_+,
\int (T_{\mathrm{test}}+\delta_n)^2\,d\nu_-
\right\}.
\end{aligned}
\]
% lean: CausalSmith.Stat.AnnotationRarearmFrontier.ateFunctional_mem_Icc
% lean: CausalSmith.Stat.AnnotationRarearmFrontier.labelFloor_bounded_mse
% lean: Causalean.Stat.half_one_sub_tvDist_le_max_error
% lean: CausalSmith.Stat.AnnotationRarearmFrontier.labelFloor_testing_mse

\item \textit{The original-record minimax bound.}
Fix any clipped Borel randomized rule \(T\) admitted by
\cref{obj:def:risk}. Apply the testing inequality with
\[
\nu_\pm:=\mathsf P_{\mathrm{ann},\pm},
\qquad
T_{\mathrm{test}}(\mathcal D,U):=T(\mathcal D,U).
\]
Using the two target identities, we obtain
\[
\max\left\{
\mathbb E_{\mathsf P_{\mathrm{ann},+}}
[(T(\mathcal D,U)-\tau(P_{+,n}))^2],
\mathbb E_{\mathsf P_{\mathrm{ann},-}}
[(T(\mathcal D,U)-\tau(P_{-,n}))^2]
\right\}
\ge\frac{\delta_n^2}{4}
=c\,b(n,\epsilon).
\]
Both laws belong to \(\mathcal M_{d,\epsilon}\), so
\[
\sup_{P\in\mathcal M_{d,\epsilon}}
\mathbb E_{P^{\otimes n}\otimes P_{XA}^{\otimes m}
                   \otimes\mathrm{Uniform}[0,1]}
[(T(\mathcal D,U)-\tau(P))^2]
\ge c\,b(n,\epsilon).
\]
The supremum is finite by the bound of four, and the rule class is
nonempty. Taking the infimum over all such \(T\), as specified in
\cref{obj:def:risk}, proves
\[
R(n,m,d,\epsilon)\ge c\,b(n,\epsilon).
\]
% lean: Causalean.Stat.le_minimaxValue_of_two_point
% lean: CausalSmith.Stat.AnnotationRarearmFrontier.rare_label_floor

\item \textit{The exact-marginal minimax bound.}
Define the complete-record laws augmented by the seed by
\[
\mathsf P_{\mathrm K,\pm}
:=
\mathsf P_{\mathrm{lab},\pm}
\otimes\mathrm{Uniform}[0,1].
\]
The common-factor identity gives
\[
\operatorname{TV}
(\mathsf P_{\mathrm K,+},\mathsf P_{\mathrm K,-})
=
\operatorname{TV}
(\mathsf P_{\mathrm{lab},+},\mathsf P_{\mathrm{lab},-})
\le\frac14.
\]
Fix a clipped Borel rule \(T^{\mathrm K}\) admitted by
\cref{obj:def:known-risk}, and freeze its marginal input at the common
table:
\[
\widetilde T^{\mathrm K}(\mathcal L,U)
:=
T^{\mathrm K}(\mathcal L,\pi_{XA},U).
\]
This is a measurable clipped rule on the complete-record arrays and the
seed. Applying the testing inequality to this rule gives
\[
\max\left\{
\mathbb E_{\mathsf P_{\mathrm K,+}}
[(\widetilde T^{\mathrm K}(\mathcal L,U)-\delta_n)^2],
\mathbb E_{\mathsf P_{\mathrm K,-}}
[(\widetilde T^{\mathrm K}(\mathcal L,U)+\delta_n)^2]
\right\}
\ge\frac{\delta_n^2}{4}
=c\,b(n,\epsilon).
\]
Because \((P_{\pm,n})_{XA}=\pi_{XA}\), these are precisely the two risks
of \(T^{\mathrm K}\) at \(P_{+,n}\) and \(P_{-,n}\). Consequently,
\[
\sup_{P\in\mathcal M_{d,\epsilon}}
\mathbb E_{P^{\otimes n}\otimes\mathrm{Uniform}[0,1]}
\left[
\bigl(T^{\mathrm K}(\mathcal L,P_{XA},U)-\tau(P)\bigr)^2
\right]
\ge c\,b(n,\epsilon).
\]
Again the supremum is bounded by four and the rule class contains the
constant zero rule. The squared-error clipping reduction in
\cref{obj:def:risk} also applies when the exact marginal is supplied as
an input: projecting the rule's output onto \([-1,1]\) preserves Borel
measurability and weakly decreases squared loss, since
\(\tau(P)\in[-1,1]\). Thus the infimum over all Borel exact-marginal
rules equals the infimum over clipped Borel exact-marginal rules.
Taking the infimum in \cref{obj:def:known-risk}
proves
\[
R^{\mathrm K}(n,d,\epsilon)\ge c\,b(n,\epsilon).
\]
The chosen \(c=1/1024>0\) is independent of all public indices.
% lean: CausalSmith.Stat.AnnotationRarearmFrontier.labelFloor_tv_common_product
% lean: Causalean.Stat.le_minimaxValue_of_two_point
% lean: CausalSmith.Stat.AnnotationRarearmFrontier.rare_label_floor

\item \textit{The pair at an arbitrary admissible separation.}
Now fix \(d\ge2\), \(0<\epsilon\le1/4\), and \(0<\delta\le1/8\).
Define the asserted pair by
\[
\begin{aligned}
P_\pm(j,1,1)&=\frac{\epsilon}{2}\left(\frac12\pm\delta\right),&
P_\pm(j,1,0)&=\frac{\epsilon}{2}\left(\frac12\mp\delta\right),\\
P_\pm(j,0,1)&=\frac{1-\epsilon}{4},&
P_\pm(j,0,0)&=\frac{1-\epsilon}{4},
&&j=1,2,
\end{aligned}
\]
and
\[
P_\pm(j,a,y)=0,
\qquad j>2,\quad a,y\in\{0,1\}.
\]
All occupied-cell entries are positive, and their sums are
\[
\sum_{a,y}P_\pm(j,a,y)
=
\frac{\epsilon}{2}+\frac{1-\epsilon}{2}
=\frac12,
\qquad j=1,2.
\]
Thus each table is a probability law. Its occupied-cell quantities are
\[
p_j(P_\pm)=\frac12,
\qquad
e_j(P_\pm)=\frac{\epsilon/2}{1/2}=\epsilon,
\qquad
\mu_{0j}(P_\pm)=\frac12,
\qquad
\mu_{1j}(P_\pm)=\frac12\pm\delta.
\]
Every other covariate mass is zero. Since
\(\epsilon\in[\epsilon,1-\epsilon]\), \cref{obj:def:model} gives
\[
P_+,P_-\in\mathcal M_{d,\epsilon}.
\]
The two occupied contributions in \cref{obj:def:target} give
\[
\tau(P_\pm)
=
\frac12(\pm\delta)+\frac12(\pm\delta)
=\pm\delta.
\]
Finally, summing over outcomes gives
\[
(P_\pm)_{XA}(j,a)=
\begin{cases}
\epsilon/2,&j\in\{1,2\},\ a=1,\\
(1-\epsilon)/2,&j\in\{1,2\},\ a=0,\\
0,&j>2,
\end{cases}
\]
which proves \((P_+)_{XA}=(P_-)_{XA}\). The null-cell assertions are
vacuous when \(d=2\).
% lean: CausalSmith.Stat.AnnotationRarearmFrontier.labelFloor_model
% lean: CausalSmith.Stat.AnnotationRarearmFrontier.labelFloor_target
% lean: CausalSmith.Stat.AnnotationRarearmFrontier.labelFloor_auxMarginal
% lean: CausalSmith.Stat.AnnotationRarearmFrontier.labelFloor_cellMass
% lean: CausalSmith.Stat.AnnotationRarearmFrontier.labelFloor_propensity
% lean: CausalSmith.Stat.AnnotationRarearmFrontier.labelFloor_outcomeMean

\item \textit{Exact divergence at separation \(1/8\).}
At \(\delta=1/8\), the treated success and failure masses in each
occupied cell under \(P_+\) are \(5\epsilon/16\) and
\(3\epsilon/16\), respectively; under \(P_-\) they are exchanged.
The control terms have likelihood ratio one, and the null cells have
zero contribution. Therefore, for each \(j=1,2\),
\[
\begin{aligned}
\sum_{a,y}
P_+(j,a,y)\log\frac{P_+(j,a,y)}{P_-(j,a,y)}
&=
\frac{5\epsilon}{16}\log\frac53
+
\frac{3\epsilon}{16}\log\frac35\\
&=
\frac{\epsilon}{8}\log\frac53,
\end{aligned}
\]
using \(\log(3/5)=-\log(5/3)\). Summing the two occupied-cell
contributions proves
\[
D(P_+\Vert P_-)
=
\sum_{z:\,P_+(z)>0}
P_+(z)\log\frac{P_+(z)}{P_-(z)}
=
\frac{\epsilon}{4}\log\frac53.
\]
% lean: CausalSmith.Stat.AnnotationRarearmFrontier.labelFloor_KL_eighth
\end{enumerate}
\end{proof}
